\section{Uniform line metrics and the split endomorphism gap}
\label{sec:line-backgrounds}

All norms in this section use the actual Ricci-flat metric $g_p$.
The positive scalar Laplacian is $\Delta=\dd^*\dd$, and
$\dd V_p=\omega_p^3/3!$. For a line metric $h$ we use
$F_h=\dbar(h^{-1}\partial h)$, so that
$c_1(L)=[\mathrm iF_h/(2\pi)]$. In particular
\begin{equation}\label{eq:line-conformal}
 \mathrm i\Lambda\partial\dbar u=-\tfrac12\Delta u,
 \qquad F_{h^{\rm ref}e^{-u}}=F_{h^{\rm ref}}+\partial\dbar u.
\end{equation}
Let $h_p^{\rm ref}$ be the restriction of the product line metric
on $\OO_{\mathbf P}(1,-1)$.

\begin{proposition}\label{prop:line-he}
There is a unique real function $u_p$ of mean zero such that
$h_p=h_p^{\rm ref}e^{-u_p}$ has $\mathrm i\Lambda F_{h_p}=0$.
The functions $u_p$ are uniformly bounded on $\mathcal B^\circ$
near zero, and have uniform bounds in every spatial derivative
on each fixed compact exterior. The curvature form
$\alpha_p=\mathrm iF_{h_p}/(2\pi)$ is the primitive harmonic
representative of $D$, with
\begin{equation}\label{eq:line-curvature-mass}
 \|\alpha_p\|_2^2=16,
 \qquad\|F_{h_p}\|_2^2=64\pi^2.
\end{equation}
\end{proposition}

\subsection{Diameter and the scalar inverse}

We first establish the geometric estimates used in the proof:
\begin{equation}\label{eq:diameter-ambient}
 \operatorname{diam}(X_p,g_p)\le D_*,\qquad
                  g_p\ge c\,\iota_p^*g_{\mathbf P}.
\end{equation}
Here $g_{\mathbf P}$ is a fixed product Fubini--Study metric.
Choose one auxiliary radius $R$ in
\eqref{eq:ricci-flat-input} so small that its zeroth-order
error is less than $1/4$. For all sufficiently small $p$,
$g_p$ is uniformly comparable with $\widetilde g_p^R$.
On a node chart the normalized Stenzel metric satisfies
\begin{equation}\label{eq:stenzel-euclidean-lower}
 \sigma_{i,p}\ge c(|z|^2+|\tau_i|)^{-1/3}
                              g_{\rm eucl}|_{Q(z)=\tau_i}.
\end{equation}
To check the normalization, the metric on $Q(Z)=1$ is bounded
below by $c(1+|Z|^2)^{-1/3}g_{\rm eucl}$, by positivity on
the compact core and the conic expansion on its end.
The change $z=\tau_i^{1/2}Z$ scales the Stenzel metric by
$|\tau_i|^{2/3}$ and the Euclidean metric by $|\tau_i|$.
This proves \eqref{eq:stenzel-euclidean-lower}; phase rotations
are unitary and the volume-normalizing constants stay in a
compact positive interval.

The embedding $\iota$ is holomorphic on the total local
hypersurface. Its coordinate functions extend holomorphically
to the ambient chart, so their spatial derivatives are uniformly
bounded on smaller charts. Hence
$\iota_p^*g_{\mathbf P}\le Cg_{\rm eucl}$ there.
On the fixed transition annuli and exterior the initial metrics
converge smoothly to a positive metric. This proves the second
inequality of \eqref{eq:diameter-ambient}. A Stenzel region
of fixed conic radius has uniformly bounded intrinsic diameter:
its core has diameter $O(|\tau_i|^{1/3})$, and its conical
end has bounded radial length and angular diameter.
The connected fixed exterior has bounded intrinsic diameter
by smooth convergence. Joining each cap to that exterior
proves the first inequality. In particular the logarithmic
length governing capacity does not produce unbounded intrinsic
diameter. The ambient reference curvature consequently satisfies
\begin{equation}\label{eq:reference-contraction}
 \|\mathrm i\Lambda_{\omega_p}F_{h_p^{\rm ref}}\|_\infty
                                                    \le C_{\rm ref}.
\end{equation}

For completeness we derive the uniform scalar estimate used
here from the Li--Yau Harnack inequality. On a connected
compact real $n$-manifold with nonnegative Ricci curvature
and diameter at most $D_*$, let $H_t(x,y)$ be the kernel of
$e^{-t\Delta}$. The inequality
\cite[Theorem 2.3]{LiYau}, with zero lower Ricci bound, gives
\[
 H_s(x,z)\le H_t(y,z)(t/s)^{n/2}
             \exp\!\left(\frac{d(x,y)^2}{4(t-s)}\right).
\]
Set $T=D_*^2$, $s=T/2$, $t=T$, and integrate in $x$.
The heat kernel has mass one, so
\[
 H_T(y,z)\ge\frac{c_n}{\Vol(X)},\qquad
                        c_n=2^{-n/2}e^{-1/2}.
\]
Thus $e^{-T\Delta}=c_n\Pi+(1-c_n)Q$, where $\Pi$ is
averaging and $Q$ has a nonnegative mass-one kernel.
For mean-zero bounded $f$, iteration and the $L^\infty$
contraction at intermediate times yield
\begin{equation}\label{eq:scalar-sup-inverse}
 \Delta^{-1}f=\int_0^\infty e^{-t\Delta}f\,\dd t,
 \qquad
 \|\Delta^{-1}f\|_\infty
        \le 2^{n/2}\sqrt e\,D_*^2\|f\|_\infty.
\end{equation}
The integral converges uniformly. Testing against smooth
functions proves the equation, and ordinary elliptic
regularity on each smooth fibre gives the classical solution.
It has mean zero and is unique. No injectivity-radius or
sectional-curvature lower bound enters this argument.

\begin{proof}[Proof of Proposition~\ref{prop:line-he}]
Take $f_p=2\mathrm i\Lambda F_{h_p^{\rm ref}}$.
By \eqref{eq:chern-numbers},
\[
 \int_{X_p} f_p\,\dd V_p
 =2\int_{X_p}\mathrm iF_{h_p^{\rm ref}}\wedge\frac{\omega_p^2}{2}
 =2\pi DH^2=0.
\]
Apply \eqref{eq:scalar-sup-inverse} in real dimension six
and use \eqref{eq:reference-contraction}. The mean-zero
solution of $\Delta_pu_p=f_p$ satisfies
\[
 \|u_p\|_\infty\le16\sqrt e\,D_*^2C_{\rm ref}.
\]
Equation \eqref{eq:line-conformal} proves the HE condition
and uniqueness. In particular $h_p$ and $h_p^{\rm ref}$ are
globally comparable. On a fixed exterior the coefficients
of the equation have uniform bounds in all spatial derivatives
and uniform ellipticity. Interior estimates on a slightly
larger exterior give the stated derivative bounds for
$u_p$ and its Chern connection. Fibrewise line-bundle
existence is also \cite[Proposition 4.2.4]{Kobayashi}; the
argument above supplies the uniform bound on this family.

The real form $\alpha_p$ is closed, of type $(1,1)$, and
primitive. On a Kähler threefold $*\alpha_p=-\alpha_p\wedge\omega_p$,
so it is coclosed and harmonic. Therefore
\[
 \|\alpha_p\|_2^2=-\int_{X_p}\alpha_p^2\wedge\omega_p
                            =-D^2H=16.
\]
This gives \eqref{eq:line-curvature-mass} with the specified
curvature convention.
\end{proof}

\subsection{The complete split kernel and gap}

Equip $E^{\rm sp}_p=L_p\otimes\CC^2\oplus L_p^{-1}\otimes\CC^2$
with $h_{0,p}=h_p\otimes I_2\oplus h_p^{-1}\otimes I_2$.
Let $D_{0,p}$ be the induced connection on endomorphisms.
It will be abbreviated to $D_0$ on a fixed fibre.
The section Bochner identity
\cite[Chapter VII, Corollary 1.3]{Demailly} gives
\begin{equation}\label{eq:split-bochner}
 D_0^*D_0=2\dbar_0^*\dbar_0,
\end{equation}
because the contracted curvature acts by commutation with
$\mathrm i\Lambda F_{h_{0,p}}=0$.

\begin{proposition}\label{prop:split-gap}
The kernel of $D_0^*D_0$ on complex endomorphism sections is
exactly the constant diagonal-block space
$\mathfrak{gl}_2\oplus\mathfrak{gl}_2$. There is $c>0$,
independent of all node-scale ratios, such that
\begin{equation}\label{eq:split-section-gap}
 \|D_0s\|_2^2\ge c\|s\|_2^2
 \quad\text{if }s\perp\ker D_0.
\end{equation}
On trace-free Hermitian sections the kernel is
\begin{equation}\label{eq:seven-space}
 \mathfrak h=\{(S_+,S_-)\in\Herm_2^2:
                             \tr S_++\tr S_-=0\},
\end{equation}
a real seven-dimensional space, with exact mass
\begin{equation}\label{eq:split-mass}
 M(S,T)=V\{\tr(S_+T_+)+\tr(S_-T_-)\}.
\end{equation}
\end{proposition}

\begin{proof}
The endomorphism bundle splits orthogonally into the two
trivial matrix blocks and the $L^{\pm2}$ matrix blocks.
The former have the constant matrices as kernel, and the
scalar Poincaré estimate controls their orthogonal complement.
The latter have no holomorphic sections, by
Proposition~\ref{prop:effective-bundles}; their covariant and
Dolbeault section energies agree up to the factor two in
\eqref{eq:split-bochner}. We must show their gaps are uniform.

Suppose, to the contrary, that $p_j\to0$ and sections of
$L_{p_j}^2$ satisfy
$\|s_j\|_2=1$ and $\|D_0s_j\|_2\to0$.
The Kato inequality and scalar Poincaré estimate imply
\begin{equation}\label{eq:kato-no-loss}
 \big\||s_j|-a_j\big\|_2\longrightarrow0,
 \qquad a_j=V^{-1}\int|s_j|\,\dd V_{p_j}
                         \longrightarrow V^{-1/2}.
\end{equation}
Indeed $|\dd|s_j||\le|D_0s_j|$ almost everywhere, with
regularization at the zero set, and
$1=\||s_j|-a_j\|_2^2+Va_j^2$.
Choose a connected fixed compact exterior $C$ containing
all topology outside small node neighbourhoods. The relative
line bundle gives fixed smooth identifications there.
Proposition~\ref{prop:line-he} allows a subsequence whose
connections and metrics converge in every spatial derivative
on a slightly larger exterior. In finitely many frames,
the sections are bounded in $W^{1,2}$. Rellich compactness
and \eqref{eq:kato-no-loss} give a smooth parallel limit
$s_0$ on the interior of $C$, with constant norm $V^{-1/2}$.
Thus the limiting connection of $L^2$ would be flat on $C$.

This contradicts the line bundle's fixed exterior class.
Removing the exceptional curves from the small resolution
$Y$ preserves $H^2(Y;\RR)$: their real codimension is four,
so the relative groups in degrees two and three in the Thom
sequence vanish. The punctured resolution retracts onto $C$.
The class of $L$ there is $D_Y=k_1-k_2$, and
$D_Y^2H=-16$, hence $2D_Y|_C\ne0$. The limiting curvature
still represents this class, either by its fixed reference
connection and scalar correction, or by integration over
fixed exterior two-cycles. It cannot vanish.
The same argument applies to $L^{-2}$.

Failure of a uniform bound on any smaller period ball would
produce the sequence just excluded, irrespective of the
relative node scales. Taking the minimum of the scalar and
off-diagonal gaps proves \eqref{eq:split-section-gap}.
The parallel Hermitian blocks each have real dimension four;
one trace condition leaves seven. Their pointwise norm is
the ordinary Hilbert--Schmidt norm, because the line metric
cancels on endomorphisms. Integrating gives
\eqref{eq:split-mass}.
\end{proof}

These estimates separate two geometric roles. The curvature
$\alpha_p$ represents the surviving divisor class $D$ and
has the fixed mass \eqref{eq:line-curvature-mass}. The weak
endomorphisms to be studied are deformations of
\eqref{eq:seven-space}. Neither space is the lost two-form
cohomology quotient governing the positive small Hodge
spectrum. The base metrics are shared, while the operators
and their finite-dimensional spaces are different.
